% ============================================================================
% Teraval — Corporate Finance individual report
% Compile on Overleaf with pdfLaTeX.
% REQUIRED IMAGE FILES (upload to the SAME Overleaf project, names must match):
%   teraval-overview.png     — Overview tab (3D hall + KPI cards + AI recommendation)
%   teraval-cashflow.png     — Cash Flow tab (FCF bars + cumulative discounted line)
%   teraval-scenarios.png    — Scenarios tab (comparison chart + table)
%   teraval-buildvsrent.png  — Build vs Rent tab (Equivalent Annual Cost)
% ============================================================================
\documentclass[12pt,a4paper]{article}

\usepackage[a4paper,top=2.6cm,bottom=2.6cm,left=2.4cm,right=2.4cm]{geometry}
\usepackage{amsmath}
\usepackage{booktabs}
\usepackage{tabularx}
\usepackage{graphicx}
\usepackage{caption}
\usepackage{enumitem}
\usepackage{parskip}
\usepackage[hidelinks]{hyperref}

% --- single-line black border on EVERY page ---
\usepackage{eso-pic}
\usepackage{tikz}
\usetikzlibrary{calc}
\AddToShipoutPictureBG{%
  \begin{tikzpicture}[remember picture,overlay]
    \draw[line width=0.6pt]
      ($(current page.north west)+(1cm,-1cm)$) rectangle
      ($(current page.south east)+(-1cm,1cm)$);
  \end{tikzpicture}%
}

\captionsetup{font=small,labelfont=bf}
\setlength{\parskip}{5pt}
\renewcommand{\arraystretch}{1.15}

% Screenshot include that degrades gracefully: if the PNG is not yet uploaded
% (they are captured manually and kept out of git), show a labelled placeholder
% instead of failing the compile. On Overleaf with the PNGs present, identical.
\newcommand{\dashfig}[2]{%
  \IfFileExists{#1.png}%
    {\includegraphics[width=#2\textwidth]{#1}}%
    {\fbox{\parbox{0.88\textwidth}{\centering\vspace{1.6cm}\small\itshape dashboard screenshot pending: #1.png\vspace{1.6cm}}}}%
}

\begin{document}

% ============================ TITLE PAGE ============================
\begin{titlepage}
\centering
\vspace*{1.2cm}
{\large\scshape Corporate Finance\par}
\vspace{0.4cm}
{\rule{0.6\textwidth}{0.4pt}\par}
\vspace{1.4cm}
{\Huge\bfseries Teraval\par}
\vspace{0.5cm}
{\Large An AI-Enabled Corporate-Finance\\[2pt] Capital-Budgeting Decision Dashboard\par}
\vspace{0.6cm}
{\large\itshape Should Barq AI build a 40\,MW AI/GPU data-centre in Abu Dhabi?\par}
\vspace{1.6cm}

{\normalsize Submitted for the partial fulfilment of the requirements of the\\
\textbf{Corporate Finance} module of the Masters in Artificial Intelligence with Business.\par}
\vspace{1.6cm}

\begin{tabular}{rl}
\textbf{Subject} & Corporate Finance\\[3pt]
\textbf{Instructor} & Dr.\ Nathaniel Christopher\\[3pt]
\textbf{Team members} & Kartik Joshi\\
 & Prem Kukreja\\
 & Gagandeep Singh\\
 & Samuel Alex\\
 & Aditya Chitale\\
\end{tabular}

\vfill
{\small Topic 9 --- AI Capital-Budgeting Dashboard (subsuming Topic 10, Relevant Cash-Flows).\par}
\end{titlepage}

% ============================ CONTENTS ============================
\tableofcontents
\thispagestyle{empty}
\newpage

% ============================ 1 ============================
\section{Executive Summary}
Teraval appraises whether \textbf{Barq AI} --- a realistic UAE AI-cloud operator, benchmarked
to the Stargate\,UAE / Khazna / G42 build-out --- should commit roughly \textbf{AED 5.8 billion}
to build and operate a \textbf{40\,MW AI/GPU data-centre hall} in Abu Dhabi over an eight-year
horizon. The base case \emph{creates value}: a net present value of \textbf{+AED 1{,}854\,m},
an internal rate of return of \textbf{16.5\%} against a 9\% cost of capital, and a profitability
index of \textbf{1.31}. The margin of safety is thin, however: value turns negative once the GPU
rental rate falls below a \textbf{break-even of about \$3.34/GPU-hr}, which sits inside today's
volatile market band of \$2.85--3.50. The recommendation is therefore to \textbf{accept, with
conditions} --- de-risk revenue through contracted offtake and stage the capital --- while keeping
a cloud-rental fallback. This report and the accompanying dashboard together answer the five
project questions: the decision and its importance (\S2), the concepts and data required (\S3--4),
the results across scenarios (\S8), how AI improves the analysis (\S6), and the final
recommendation with its risks (\S9--10).

% ============================ 2 ============================
\section{Financial Problem}
The decision-maker is the CFO and investment committee of Barq AI. The problem is a
near-irreversible, ``bet-the-company'' capital commitment: about AED 5.8 billion of up-front
capital sunk into an asset whose revenue --- the price of renting GPU compute --- \emph{collapsed
from roughly \$8 to \$2.85--3.50 per GPU-hour} across 2024--25 as supply flooded the market. The
decision matters because of the sheer scale of capital at risk, the volatility of the revenue
line, and the strategic stakes of the UAE's ambition to become a global AI-compute hub. Getting
it wrong means either a multi-billion-dirham value destruction or forgoing a defining
infrastructure position. For scale, the reference Stargate\,UAE campus commits roughly
\$10 billion for a 1\,GW build, of which a 200\,MW first phase completes in 2026; Barq's 40\,MW
hall is one realistic module of exactly this kind of investment. The capital is effectively
irreversible --- purpose-built power and liquid cooling cannot be redeployed --- so the
opportunity cost of a wrong decision is close to the entire outlay. This is the classic
capital-budgeting accept/reject question, and the natural capstone of the course's
project-appraisal and capital-investment units.

% ============================ 3 ============================
\section{Relevant Corporate Finance Concepts}
The appraisal rests on the time value of money and discounted cash flow. Only \emph{relevant
(incremental) cash flows} enter the model --- the initial outlay, incremental operating cash
flows, and the terminal value --- while sunk costs are excluded and working capital and salvage
are included. The core metrics are:
\begin{align*}
\text{NPV} &= \sum_{t=1}^{n}\frac{CF_t}{(1+r)^t}-CF_0, &
\text{PI} &= \frac{\text{PV of future flows}}{|CF_0|},\\
\text{IRR: } 0 &= \sum_{t=0}^{n}\frac{CF_t}{(1+\text{IRR})^t}, &
\text{EAC} &= \frac{\text{PV of costs}}{\big(1-(1+r)^{-n}\big)/r}.
\end{align*}
NPV is the primary rule (accept if positive); IRR and the profitability index confirm it; payback
and discounted payback give a liquidity view; and \emph{Equivalent Annual Cost} (EAC) is used to
compare the build against renting, because the two routes have different lives and cannot be
compared on total cost or NPV. The discount rate is the weighted average cost of capital (WACC),
with the cost of equity from the Capital Asset Pricing Model,
$K_e = R_f + \beta(R_m-R_f)$. Modified IRR is reported alongside IRR because IRR implicitly
assumes interim cash flows are reinvested at the IRR itself, which flatters a high-return project;
MIRR reinvests them at the more realistic cost of capital. Payback and discounted payback ignore
cash flows beyond the cut-off --- and, for the undiscounted version, the time value of money ---
so they serve only as a secondary liquidity check.

% ============================ 4 ============================
\section{Data and Assumptions}
Inputs are classified as historical (H), current market (C), forecast (F), user-entered (U) or
AI-generated (AI). Monetary values are AED millions; USD is converted at the pegged rate
USD\,1 = AED\,3.6725.

\begin{table}[h]
\small
\begin{tabularx}{\textwidth}{@{}l p{4.3cm} X@{}}
\toprule
\textbf{Assumption} & \textbf{Base value} & \textbf{Type / basis}\\
\midrule
GPU rental price & \$4.0 / GPU-hr & C --- cloud GPU market\\
Utilization & 80\% & U --- contracted + on-demand mix\\
Facility scale & 40\,MW IT; 300 GB200 racks (21{,}600 GPUs) & C --- GB200 NVL72 platform\\
Total capex & AED 5{,}766\,m (\$39.2\,m/MW) & C --- AI-grade build + GPUs\\
Electricity tariff & AED 0.15 / kWh & C --- Abu Dhabi industrial (official pending)\\
PUE (cooling) & 1.20 & C --- liquid-cooled GB200\\
WACC & 9\% & F --- anchored to CBUAE base 3.65\% / 3-mo EIBOR 3.94\% + risk premia\\
UAE corporate tax & 9\% & C --- introduced 2023\\
Project life & 8 years & U --- appraisal horizon\\
\bottomrule
\end{tabularx}
\caption{Key assumptions and their data type. Interest-rate anchors are official Central Bank of
the UAE figures (July 2026); the industrial tariff awaits the Abu Dhabi Open Data portal.}
\end{table}

% ============================ 5 ============================
\section{Financial Calculations}
The Year-0 outlay is AED 5{,}766\,m of capital (IT hardware AED 3{,}856\,m plus facility
AED 1{,}910\,m) plus AED 150\,m of working capital. Base-case revenue is AED 2{,}224\,m per year
against operating costs of AED 600\,m, giving EBITDA of AED 1{,}624\,m; a mid-life GPU refresh in
Year 4 and terminal salvage complete the stream. Discounting at 9\% yields the results in
Table~\ref{tab:metrics}. A minimum of three major calculations is comfortably exceeded --- NPV,
IRR, MIRR, PI, payback and break-even are all computed and cross-checked by a live model
self-test (e.g.\ NPV discounted at the IRR equals zero). The 9\% discount rate is built up from
the official Central Bank rates as a risk-free anchor, a levered equity risk premium and an
after-tax cost of debt, while straight-line depreciation supplies a tax shield inside the
operating cash flows.

\begin{table}[h]
\centering\small
\begin{tabular}{@{}lll@{}}
\toprule
\textbf{Metric} & \textbf{Base case} & \textbf{Decision signal}\\
\midrule
Net present value & +AED 1{,}854\,m & Positive $\rightarrow$ accept\\
Internal rate of return & 16.5\% & Above 9\% WACC $\rightarrow$ accept\\
Modified IRR & 12.6\% & Above WACC $\rightarrow$ accept\\
Profitability index & 1.31 & $>1 \rightarrow$ accept\\
Payback / discounted & 5.0 / 6.3 yrs & Within the 8-year life\\
Break-even GPU rate & \$3.34 / GPU-hr & Market \$2.85--3.50 $\rightarrow$ knife-edge\\
\bottomrule
\end{tabular}
\caption{Verified base-case decision metrics.}
\label{tab:metrics}
\end{table}

The model also frames the \emph{relevant} question --- build versus rent. Because Barq needs the
compute either way, the build is appraised incrementally against renting hyperscaler capacity,
compared on Equivalent Annual Cost. Building is a largely fixed cost while renting scales with
utilization, so \textbf{building wins only above roughly 77\% sustained utilization}; at the base
case it beats renting by a thin incremental NPV of about +AED 320\,m (Figure~\ref{fig:bvr}). As an
internal check, the sign of that incremental NPV agrees with the difference in equivalent annual
costs, so two independent routes to the answer are mutually consistent.

\begin{figure}[h]
\centering
\dashfig{teraval-cashflow}{0.92}
\caption{Eight-year free cash flow (bars; Year-0 outlay and the Year-4 refresh dip are negative)
and the cumulative \emph{discounted} cash flow (line), whose zero-crossing is the discounted
payback of about 6.3 years.}
\end{figure}

\begin{figure}[h]
\centering
\dashfig{teraval-buildvsrent}{0.92}
\caption{Build vs Rent on Equivalent Annual Cost. Building's fixed cost beats flexible renting
only at high, sustained utilization; below the $\sim$77\% crossover, renting is cheaper.}
\label{fig:bvr}
\end{figure}

% ============================ 6 ============================
\section{AI Features}
Teraval proposes six AI features, each grounded in the deterministic engine so that figures are
never invented:
\begin{enumerate}[leftmargin=1.4em,itemsep=1pt]
\item \textbf{Rate \& utilization forecaster} --- projects a forward GPU-rental path with
P10--P90 bands, replacing a flat guess with a defensible range (risk: market non-stationarity).
\item \textbf{Scenario generator} --- auto-builds coherent optimistic/base/pessimistic driver sets.
\item \textbf{Monte-Carlo risk simulator} --- 5{,}000 NPV draws yield a probability of loss
($\approx$21\% at base) and a value-at-risk view.
\item \textbf{AI Finance Assistant (NVIDIA NIM)} --- a natural-language, tab-aware chatbot that
answers from the live model state and never free-generates numbers (risk: hallucination, mitigated
by grounding).
\item \textbf{Anomaly / unrealistic-assumption detector} --- flags inputs outside credible ranges.
\item \textbf{Recommendation engine} --- synthesises metrics, break-even and the dominant driver
into an accept/reject conclusion.
\end{enumerate}
Together these automate the arithmetic, forecast the volatile revenue line, quantify risk, and
explain results in plain language --- directly improving the analysis, dashboard and
decision-making process.

% ============================ 7 ============================
\section{Dashboard Explanation}
The dashboard (Figure~\ref{fig:overview}) delivers the six required components across seven tabs:
KPI cards (NPV, IRR, MIRR, PI, payback, break-even); a cash-flow trend chart; a scenario
comparison; live sensitivity with a tornado chart; a risk-and-alert panel; and an AI
recommendation panel. It is a dark, glassmorphism React application whose hero is a WebGL 3D
data-centre hall that reacts to the live assumptions, with a floating, context-aware AI assistant
available on every tab, a live model-integrity self-test, and a data-provenance panel citing the
official Central Bank rates. Every component reads a single shared model state, so moving any
assumption recomputes the KPI cards, charts, scenarios, risk alerts and the recommendation
simultaneously --- letting a manager see the physical and financial story change together.

\smallskip
\noindent\textbf{Live deployment.} The dashboard is deployed publicly --- the React frontend on
Vercel at \href{https://teraval-ai.vercel.app}{teraval-ai.vercel.app}, backed by the FastAPI +
NVIDIA~NIM assistant on Render at
\href{https://teraval-assistant.onrender.com}{teraval-assistant.onrender.com}; full source at
\href{https://github.com/KartikJoshi23/TeravalAI}{github.com/KartikJoshi23/TeravalAI}. If the
backend is offline the assistant falls back to a grounded local answerer, so every quoted figure
stays reproducible from the deterministic engine.

\begin{figure}[h]
\centering
\dashfig{teraval-overview}{0.95}
\caption{The Teraval Overview: the 3D data-centre hall, the headline KPI cards, and the
AI-generated recommendation.}
\label{fig:overview}
\end{figure}

% ============================ 8 ============================
\section{Scenario and Sensitivity Analysis}
Three scenarios are modelled (Table~\ref{tab:scen}, Figure~\ref{fig:scen}). The decision
\emph{flips}: value is strongly positive in the optimistic case, thinly positive in the base case,
and sharply negative in the pessimistic case. One-way sensitivity shows the \textbf{GPU rental
price} is by far the dominant driver --- a 20\% fall alone turns NPV negative --- followed by
utilization; the electricity tariff barely moves the result, confirming that Abu Dhabi's cheap
power is a genuine competitive moat. The Monte-Carlo simulation puts the probability of a negative
NPV at about 21\% at base, with a mean simulated NPV of roughly +AED 1{,}456\,m --- below the base
case, reflecting that the downside outweighs the upside. The decision changes once the rental rate
falls below roughly \$3.34/GPU-hr or utilization drops under about 67\%; the principal financial
risk is a structural, market-wide decline in the rental rate; and the best management response is
to secure multi-year contracted offtake before committing capital, stage the fit-out, and retain a
cloud-rental fallback.

\begin{table}[h]
\centering\small
\begin{tabular}{@{}lllll@{}}
\toprule
\textbf{Scenario} & \textbf{NPV} & \textbf{IRR} & \textbf{PI} & \textbf{Decision}\\
\midrule
Optimistic & +AED 10{,}216\,m & 44.5\% & 2.73 & Accept\\
Base case & +AED 1{,}854\,m & 16.5\% & 1.31 & Accept (conditional)\\
Pessimistic & $-$AED 4{,}138\,m & $-$9.5\% & 0.30 & Reject\\
\bottomrule
\end{tabular}
\caption{Scenario results --- the accept/reject decision flips with the assumptions.}
\label{tab:scen}
\end{table}

\begin{figure}[h]
\centering
\dashfig{teraval-scenarios}{0.92}
\caption{Optimistic / base / pessimistic comparison against the live ``current'' case.}
\label{fig:scen}
\end{figure}

% ============================ 9 ============================
\section{AI Limitations and Ethical Risks}
The AI features carry real risks. Rate and utilization \emph{forecasts} are extrapolations of a
volatile, immature market and could mislead if read as certainty --- so they are shown as ranges,
not points. A language model can \emph{hallucinate} figures; the assistant is therefore grounded
in the deterministic engine and instructed to narrate only. Incorrect input data can silently flip
the verdict, so assumptions are transparent, source-tagged and range-checked. A bull-market
\emph{bias} could inflate revenue, which the pessimistic scenario and Monte-Carlo explicitly
stress. Commercially sensitive assumptions stay client-side for \emph{confidentiality}. Above all,
the tool requires \emph{human review}: it advises, but the responsibility for the final financial
decision rests with the CFO and the investment committee, not with the AI.

% ============================ 10 ============================
\section{Final Recommendation}
The recommendation is to \textbf{proceed with the investment, conditionally}. The base case
creates value (NPV +AED 1{,}854\,m; IRR 16.5\% against a conservative 9\% hurdle anchored to
official CBUAE rates), but the margin of safety is thin and the outcome hinges on one variable.
Management should therefore (i) secure multi-year contracted offtake to hold the effective rental
rate above the \$3.34 break-even; (ii) stage the GPU fit-out rather than release all capital at
once; and (iii) keep hyperscaler rental available as a fallback, since building only pays above
about 77\% sustained utilization. The dominant \emph{financial} risk is a structural collapse in
the GPU rental rate; the dominant \emph{AI-related} risk is over-reliance on optimistic forecasts.
Managed on these terms, the project is a defensible, value-creating commitment.

\end{document}
